%
\documentclass[11pt,a4paper]{moderncv}

\usepackage{adjustbox}
\usepackage{makecell}
\usepackage{tabu}
\setlength{\tabcolsep}{2pt}

\usepackage[square,comma,numbers]{natbib}
\usepackage{booktabs}
\moderncvtheme[orange]{casual}  % optional argument are 'blue' (default), 'orange', 'green', 'red', 'purple', 'grey' 
\usepackage[utf8]{inputenc} % replace by the encoding you are using
\usepackage[scale=0.80]{geometry}
\usepackage{lmodern}
\usepackage{fontawesome5}
\firstname{Thibault}
\familyname{Latrille}
\usepackage[dvipsnames]{xcolor}
\definecolor{BLUELINK}{HTML}{0645AD}
\definecolor{DARKBLUELINK}{HTML}{0B0080}
\definecolor{LIGHTGREY}{gray}{0.9}
\PassOptionsToPackage{hyphens}{url}
\usepackage[colorlinks=false]{hyperref}
% for linking between references, figures, TOC, etc in the pdf document
\hypersetup{colorlinks,
linkcolor=DARKBLUELINK,
anchorcolor=DARKBLUELINK,
citecolor=DARKBLUELINK,
filecolor=DARKBLUELINK,
menucolor=DARKBLUELINK,
urlcolor=BLUELINK
} % Color citation links in purple
\PassOptionsToPackage{unicode}{hyperref}
\PassOptionsToPackage{naturalnames}{hyperref}


\renewcommand{\cellalign}{cl}
\renewcommand{\theadalign}{cl}

\begin{document}

\maketitle
\vspace*{-1cm}
\cventry{\faInternetExplorer}{\href{https://thibaultlatrille.github.io}{thibaultlatrille.github.io}}{}{}{}{}
\cventry{\faOrcid}{\href{https://orcid.org/0000-0002-9643-4668}{orcid.org/0000-0002-9643-4668}}{}{}{}{}
\cventry{\faGithub}{\href{https://github.com/ThibaultLatrille}{github.com/ThibaultLatrille}}{}{}{}{}
\cventry{\faGoogle}{\href{https://scholar.google.com/citations?user=6HlrDNgAAAAJ}{scholar.google.com/citations?user=6HlrDNgAAAAJ}}{}{}{}{}



\section{Curriculum Vitae}
\vspace*{0.5em}
{\large \textbf{Employments}}
\vspace*{0.5em}\\~
\cventry{Sep.~2021 - Aug.~2026}{Postdoctoral researcher}{Université de Lausanne, Department of Computational Biology}{advisor Nicolas Salamin.}{}{}
\cventry{Sep.~2020 - Aug.~2021}{Temporary Research and Teaching Attaché}{Université de Lyon, Biometry and Evolutionary Biology laboratory (LBBE).}{}{}{}
\cventry{Sep.~2017 - Aug.~2020}{PhD}{Université de Lyon, Biometry and Evolutionary Biology laboratory (LBBE)}{Modelling the interplay between selective and neutral mechanisms in the evolution of protein-coding DNA sequences}{under the supervision of Nicolas Lartillot and defended on November 30th, 2020.}{\href{https://tel.archives-ouvertes.fr/tel-03405159/}{tel.archives-ouvertes.fr/tel-03405159/}}
\cventry{Sep.~2011 - Aug.~2017}{Civil servant}{Ecole Normale Supérieure de Lyon (ENS)}{}{As a "élève normalien", I was a civil servant for the French state, and paid a salary by the ENS during my studies (4 years of salary). I also had to commit to work for a public institution for a duration of 10 years after signing the contract in September 2011. I have fulfilled this commitment and releaseed from this engagement in 2023 by working in public research institutions (ENS, University of Lyon and University of Lausanne) during more than 10 years.}{}

\vspace*{0.5em}
{\large \textbf{Education}}
\vspace*{0.5em}\\
\cventry{2016-2017}{Master 2 in complex systems}{Ecole Normale Supérieure de Lyon (ENS), Lyon.}{}{}{}
\cventry{2013-2014}{Master 1 in biology}{Ecole Normale Supérieure de Lyon (ENS), Lyon.}{}{}{}
\cventry{2012-2013}{Erasmus in mathematics}{University of Uppsala, Sweden.}{}{}{}
\cventry{2011-2012}{Bachelor in biology}{Ecole Normale Supérieure de Lyon (ENS), Lyon.}{}{}{}
\cventry{2009-2011}{Classes Préparatoires aux Grandes Ecoles (CPGE)}{Lycée Joffre, Montpellier.}{}{}{Post-secondary preparatory classes in Biology, Physics, Chemistry and Earth Sciences (BCPST).}

\newpage

\vspace*{0.5em}
{\large \textbf{Personal information}}
\vspace*{0.5em}\\
\cventry{Children}{One child}{Born June 11, 2023.}{}{}{}
\cvlanguage{French}{Native tongue.}{}
\cvlanguage{English}{Fluent.}{}
\cvlanguage{Spanish}{Conversational.}{}
\cventry{\faTerminal}{Python \& C++}{}{proficient use and extensive schooling.}{}{}
\cventry{\faTerminal}{Snakemake}{}{proficient use and teaching.}{}{}
\cventry{\faTerminal}{Rust, R, Scilab \& Matlab}{}{understand and have written small programs.}{}{}
\cventry{\faCode}{Web development}{}{front and back-end development.}{}{Developped website \href{https://controversciences.github.io/}{ControverSciences.org},  \href{https://symbiotager.github.io/}{MonPotager} and \href{https://wheretopublish.github.io/}{WhereToPublish}.}
\cventry{\faCodeBranch }{Versionning}{}{experimented user with git on collaborative projects.}{}{}
\cventry{\faRobot }{Machine learning}{}{Pytorch for training small ($\leq$ 32Gb) models.}{}{}
\cventry{\faPhotoVideo }{Adobe creative suite, Affinity suite}{}{experimented user.}{}{}
\cventry{\faFile*}{Office, \LaTeX.}{}{}{}{}
\cventry{\faDesktop}{Ubuntu \faUbuntu, MacOS \faApple \ and Windows \faWindows.}{}{}{}{}

\newpage

\section{Teaching}
I have extensive experience in various teaching classes (from bachelor to master level) and formats (lectures, practical sessions, workshops, and online courses) in the field of evolutionary biology, bioinformatics, and programming. I have also been involved in the design of some of these courses and workshops.\vspace*{1em}\\

\cventry{2025-2026, 1~video}{Artificial Intelligence}{lecturer}{Video for a Massive Open Online Courses (MOOC) on generative artificial intelligence for bachelor students}{University of Lausanne.}{}
\cventry{2024,\\ 1$\times$4h}{Evolution}{lecturer}{lectures on population genetics for bachelor students}{University of Lausanne.}{}
\cventry{2021-2026, 3$\times$4h}{Reproducible Science with Snakemake}{organizer}{workshop for PhD students}{University of Lausanne.}{}
\cventry{2021-2026, 5$\times$14h}{Phylogeny and Comparative Methods}{teaching assistant}{practical sessions for master students}{University of Lausanne.}{}
\cventry{2021-2026, 5$\times$12h}{Advanced Python Programming}{teaching assistant}{practical sessions for master students}{University of Lausanne.}{}
\cventry{2024-2026, 2$\times$24h}{Introduction to Programming, Bsc~1}{teaching assistant}{practical sessions for bachelor students (first year)}{University of Lausanne.}{}
\cventry{2021-2026, 5$\times$12h}{Introduction to Programming, Bsc~2}{teaching assistant}{practical sessions for bachelor students (second year)}{University of Lausanne.}{}
\cventry{2024,\\ 1$\times$32h}{Modeling in Biology and Bioinformatics}{teaching assistant}{practical sessions for master students}{University of Lausanne.}{}
\cventry{2021 \& 2022, 2$\times$3h}{Scientific Methodology}{lecturer}{lectures and practical sessions for master students}{University of Montpellier.}{}
\cventry{2020-2021, 1$\times$24h}{Bioinformatics}{teaching assistant}{practical sessions for master students}{University of Lyon.}{}
\cventry{2020-2021, 1$\times$22h}{Genetics}{teaching assistant}{practical sessions for bachelor students}{University of Lyon.}{}
\cventry{2020-2021, 1$\times$18h}{Population Genetics}{teaching assistant}{practical sessions for bachelor students}{University of Lyon.}{}
\cventry{2018 \& 2019, 2$\times$12h}{Bioinformatics for Sequencing Data Analysis}{organizer}{course and practical sessions for professionals}{Laboratory of Biometry and Evolutionary Biology (LBBE).}{}
\cventry{2017-2021, 4$\times$32h}{Bioinformatics \& Statistics}{teaching assistant}{practical sessions for bachelor students}{University of Lyon.}{}
\cventry{2017-2021, 4$\times$20h}{Next-Generation Sequencing}{teaching assistant}{reproducible science for master students}{École Normale Supérieure de Lyon (ENS).}{}
\cventry{2017-2019, 2$\times$18h}{Library Research}{teaching assistant}{practical sessions for bachelor students}{University of Lyon.}{}

\newpage

\section{Projects and responsibilities}
I have been involved in various projects and responsibilities, including research projects, advising master students, editing special issues, organizing workshops, and developing websites and applications. \vspace*{1em}\\

\cventry{2025, 6~months}{Advisor for master student}{Léo Besançon}{}{}{Are signatures of selection consistent across different evolutionary scales?}
\cventry{2021-2026}{Workshop organizer}{Faculty of Biology and Medicine, Lausanne.}{}{}{I am currently organizing several workshops on reproducible science at the university, \href{https://github.com/ThibaultLatrille/workshop-snakemake-unil2025}{github.com/ThibaultLatrille/workshop-snakemake-unil2025}.}
\cventry{2025-2026}{Co-creator, Developer, and Designer}{}{\href{https://wheretopublish.github.io/}{WhereToPublish}}{}{WhereToPublish is a website providing information on the ownership of scientific journals and publishers across different fields of biology.}
\cventry{2021-2026}{Project funding}{Department of Computational Biology, Lausanne.}{}{}{I participated in writing a grant proposal for the Swiss National Science Foundation (SNSF) on evolutionary modelling across scales which was awarded to Nicolas Salamin \href{https://data.snf.ch/grants/grant/219757}{data.snf.ch/grants/grant/219757}. This grant is currently funding my salary, but it could only be deposited by the principal investigator under his sole name (Nicolas Salamin).}
\cventry{2023-2024}{Guest Editor}{Journal of Evolutionary Biology.}{}{}{Editor on a Special Issue on the integration of micro- and macro-evolution, five of the manuscripts that have handled as a Guest Editor have now been published (December 2024 volume of the journal, \href{https://academic.oup.com/jeb/issue/37/12}{academic.oup.com/jeb/issue/37/12}).}
\cventry{2023, 6~months}{Advisor for master student}{Theo Schneider}{}{}{Are fitness landscapes stable through time?}
\cventry{2015-2026}{Creator, Developer, and Designer}{}{\href{https://controversciences.github.io/}{ControverSciences.org}}{} {ControverSciences.org is a website where contributors collaboratively publish summaries of scientific controversies, based on and analyzing scientific publications. The site has been used as a teaching tool in the Master's program at the University of Montpellier, in a course led by Catherine Moulia.}
\cventry{2020-2026}{Creator, Developer, and Designer}{}{\href{https://symbiotager.github.io/}{MonPotager}}{} {MonPotager is an application that simulates a vegetable garden by allowing users to insert various species of fruits and vegetables, and determine whether their interactions will be favorable or unfavorable. The site is used as a teaching tool in the Bachelor's program at INSA Lyon, in a course led by Christophe Rigotti.}
\cventry{2019}{Project funding}{Université de Lyon.}{}{}{Participated in designing and writing a grant proposal (Community Garden Book) for Action Exploratoire (AEx) that was awarded to Eric Tannier.}
\cventry{2012}{Laboratoire Junior}{Ecole Normale Supérieure de Lyon (ENS).}{}{}{ENS promotes and helps projects built by students at their own initiative. Our project was to write a software able to produce 3D map of the drosophila's composed eye, based on pictures that cover a small part of the eye (picture matching).}


\newpage
\section{Collective engagements}
I have been involved in various collective engagements, including being part of PhD committees, organizing conferences and symposia, and participating in the organization of department retreats, as well as being part of an association to maintain and improve scientific animation and cohesion activities.\vspace*{1em}\\

\cventry{2025-2026}{PhD committee}{Université de Montpellier.}{}{}{I have been part of the PhD committee of Emeline Esnouf.}
\cventry{2025}{Symposium organization}{ESEB 2025, Barcelona.}{}{}{Co-organizer of the symposium \href{https://eseb2025.com/list-of-symposia/}{Microevolutionary processes and Macroevolutionary patterns (S36)}.}
\cventry{2024-2025}{Conference organization}{Department of Ecology and Evolution, Lausanne.}{}{}{Part of the scientific committee for the conference \href{https://wp.unil.ch/biology25/}{Biology 2025}, held February 13-14, 2025 at the University of Lausanne.}
\cventry{2022-2026}{Department retreat organization}{Department of Computational Biology, Lausanne.}{}{}{I have been part of the assistant committee. I also organized a workshop to strengthen connection between labs, \href{https://thibaultlatrille.github.io/DBCNet}{thibaultlatrille.github.io/DBCNet}. I organized a workshops to collaboratively write a \href{https://wiki.unil.ch/fbm-dbc/}{internal wiki} for the department.}
\cventry{2021-2024}{PhD committee}{Laboratoire de Biométrie et Biologie Évolutive (LBBE), Lyon.}{}{}{I have been part of the PhD committee of Melodie Bastian, who defended her thesis in 2024.}
\cventry{2019-2020}{President}{Les pinsons migRateurs.}{}{}{I have co-founded the association and was president for its first term, the association goal is to maintain and improve scientific animation, cohesion activities and participate in professional integration within the Laboratoire de Biométrie et Biologie Évolutive (LBBE). We created "Le guide de doctorants", an informal guide to help navigate administrative mazes and share our best practices as researchers and teachers.}
\cventry{2017-2020}{Expert.}{}{}{}{Participation in the AREN (\href{https://e-fran.education.gouv.fr/aren/}{e-fran.education.gouv.fr/aren}) project: ARgumentation et Numérique  as a provider of content to be debated by high school students (\href{https://github.com/ThibaultLatrille/AREN-corpus}{github.com/ThibaultLatrille/AREN-corpus}).}
\cventry{2015-2020}{Jury}{One week per year.}{}{}{Expert as founder of the ControverSciences.org platform used in the teaching of 'Scientific controversies' in master 2 at the University of Montpellier in a course led by Catherine Moulia.}

\newpage
\section{Publication list}
Publication List which includes (i) peer-reviewed articles, (ii) preprints and articles without peer review, (iii) invited talks, and (iv) conference talks.
Publications are ordered from most recent to oldest, with full references, working links, and clear indication of co-first authors ($^\dag$) and PhD advisor (\textcolor{Salmon}{N. Lartillot}).

\section*{(i) Peer-reviewed publications}
\begin{enumerate}
\setlength\itemsep{1em}
    \item \textbf{T. Latrille}, T. Gaboriau, and N. Salamin. Phylograms better fit neutrally evolving traits than chronograms. \textit{Evolution Letters}, qrag015, Apr. 2026. \href{https://doi.org/10.1093/evlett/qrag015}{doi.org/10.1093/evlett/qrag015}.
    \item L. M. Fitzgerald, \textbf{T. Latrille}, A. Marcionetti, T. Gaboriau, D. A. Hartasánchez, and N. Salamin. Genetic basis of the evolution of vertical bars in clownfishes. \textit{BMC genomics}, 27(216), Jan. 2026. \href{https://doi.org/10.1186/s12864-026-12570-9}{doi.org/10.1186/s12864-026-12570-9}.
    \item E. Trucchi, P. Massa, F. Giannelli, \textbf{T. Latrille}, M. Gargano, F. A. N. Fernandes, L. Ancona, N. C. Stenseth, J. F. Obiol, J. Paris, G. Bertorelle, and C. L. Bohec. High gene expression predicts extremely low segregation of deleterious mutations in large penguin populations. \textit{Molecular Biology and Evolution}, 42(6):msaf146, 2025. \href{https://doi.org/10.1093/molbev/msaf146}{doi.org/10.1093/molbev/msaf146}.
    \item \textbf{T. Latrille}$^{\dag}$, J. Joseph$^{\dag}$, D. A. Hartasánchez, and N. Salamin. Estimating the proportion of beneficial mutations that are not adaptive in mammals. \textit{PLOS Genetics}, 20(12):e1011536, Dec. 2024. \href{https://doi.org/10.1371/journal.pgen.1011536}{doi.org/10.1371/journal.pgen.1011536}.
    \item M. Tsuboi, T. Gaboriau, and \textbf{T. Latrille}. An introduction to the special issue: inferring macroevolutionary patterns and processes from microevolutionary mechanisms \textit{Journal of Evolutionary Biology}, 37(12):1395–1401, Dec. 2024. \href{https://doi.org/10.1093/jeb/voae132}{doi.org/10.1093/jeb/voae132}.
    \item \textbf{T. Latrille}, M. Bastian, T. Gaboriau, and N. Salamin. Detecting diversifying selection for a trait from within and between-species genotypes and phenotypes. \textit{Journal of Evolutionary Biology}, 37(12):1538–1550, Dec. 2024. \href{https://doi.org/10.1093/jeb/voae084}{doi.org/10.1093/jeb/voae084}.
    \item D. Silvestro, \textbf{T. Latrille}, and N. Salamin. Toward a Semi-Supervised Learning Approach to Phylogenetic Estimation. \textit{Systematic Biology}, page syae029, June 2024. \href{https://doi.org/10.1093/sysbio/syae029}{doi.org/10.1093/sysbio/syae029}.
    \item \textbf{T. Latrille}, N. Rodrigue, and \textcolor{Salmon}{N. Lartillot}. Genes and sites under adaptation at the phylogenetic scale also exhibit adaptation at the population-genetic scale. \textit{Proceedings of the National Academy of Sciences of the United States of America}, 120(11):e2214977120, 2023. \href{https://doi.org/10.1073/pnas.2214977120}{doi.org/10.1073/pnas.2214977120}.
    \item D. A. Hartasánchez$^{\dag}$, \textbf{T. Latrille}$^{\dag}$, M. Brasó-Vives, and A. Navarro. Bridging Time Scales in Evolutionary Biology. \textit{Mathematics Online First Collections}, pages 1–23. Springer International Publishing, Cham, 2022. \href{https://doi.org/10.1007/16618_2022_37}{doi.org/10.1007/16618\_2022\_37}.
    \item \textbf{T. Latrille} and \textcolor{Salmon}{N. Lartillot}. An Improved Codon Modeling Approach for Accurate Estimation of the Mutation Bias. \textit{Molecular Biology and Evolution}, 39(2):msac005, Feb. 2022. \href{https://doi.org/10.1093/molbev/msac005}{doi.org/10.1093/molbev/msac005}.
    \item \textbf{T. Latrille}, V. Lanore, and \textcolor{Salmon}{N. Lartillot}. Inferring long-term effective population size with mutation–selection models. \textit{Molecular Biology and Evolution}, 38(10):4573–4587, Oct. 2021. \href{https://doi.org/10.1093/molbev/msab160}{doi.org/10.1093/molbev/msab160}.
    \item \textbf{T. Latrille} and \textcolor{Salmon}{N. Lartillot}. Quantifying the impact of changes in effective population size and expression level on the rate of coding sequence evolution. \textit{Theoretical Population Biology}, 142:57–66, Dec. 2021. \href{https://doi.org/10.1016/j.tpb.2021.09.005}{doi.org/10.1016/j.tpb.2021.09.005}.
    \item N. Rodrigue, \textbf{T. Latrille}, and \textcolor{Salmon}{N. Lartillot}. A Bayesian mutation-selection framework for detecting site-specific adaptive evolution in protein-coding genes. \textit{Molecular Biology and Evolution}, 38(3):1199–1208, Mar. 2021. \href{https://doi.org/10.1093/molbev/msaa265}{doi.org/10.1093/molbev/msaa265}.
    \item \textbf{T. Latrille}. Modelling the Articulation of Selective and Neutral Mechanisms in the Evolution of Protein-Coding DNA Sequences. PhD thesis, \textit{Université de Lyon}, Nov. 2020. \href{https://tel.archives-ouvertes.fr/tel-03405159}{tel.archives-ouvertes.fr/tel-03405159}.
    \item \textbf{T. Latrille}, L. Duret, and \textcolor{Salmon}{N. Lartillot}. The Red Queen model of recombination hot-spot evolution: A theoretical investigation. \textit{Philosophical Transactions of the Royal Society of London. Series B, Biological Sciences}, 372(1736):20160463, Dec. 2017. \href{https://doi.org/10.1098/rstb.2016.0463}{doi.org/10.1098/rstb.2016.0463}.
\end{enumerate}
\vspace{1em}
$^{\dag}$ Authors contributed equally to this work.\\
\textcolor{Salmon}{N. Lartillot} is the PhD advisor.

\section*{(ii) Publications without peer review}
\begin{enumerate}
\setlength\itemsep{1em}
    \item A. Laverré, \textbf{T. Latrille}, and M. Robinson-Rechavi. RegEvol: detection of directional selection in regulatory sequences through phenotypic predictions and phenotype-to-fitness functions. \textit{bioRxiv}, 2025.11.26.690685, Nov. 2025. \href{https://doi.org/10.1101/2025.11.26.690685}{doi.org/10.1101/2025.11.26.690685}.
    \item B. M. Farina, \textbf{T. Latrille}, N. Salamin, D. Silvestro, and S. Faurby. Widespread selection relaxation in aquatic mammals, Sept. 2024. \href{https://doi.org/10.1101/2024.09.18.613479}{doi.org/10.1101/2024.09.18.613479}.
\end{enumerate}

\section*{(iii) Invited talks}
\cventry{October 3, 2024}{External Seminar, Department of Ecology and Evolution}{Lausanne.}{}{}{\href{https://ecoevo.social/@dee_unil/113204416662399007}{ecoevo.social/@dee\_unil/113204416662399007}\\Predicting Selection on Traits and Sequences: Contrast Across Evolutionary Scales.}
\cventry{April 17-21, 2023}{Biological Evolution Across Scales}{Lausanne.}{}{}{\href{https://bevas-epfl.github.io/}{bevas-epfl.github.io}\\A phylogenetic mutation-selection model predicts fitness effects of mutations in extant mammals.}

\section*{(iv) Conference talks}
\cventry{September 18–19, 2025}{Computational Biology Symposium}{Lausanne.}{}{}{\href{https://cbiosymposium.unil.ch/}{cbiosymposium.unil.ch} \\ Between and within species variation to detect selection on gene expression level in mammals.}
\cventry{August 17–22, 2025}{European Society for Evolutionary Biology}{Barcelona.}{}{}{\href{https://eseb.org/wp-content/uploads/2025/09/FULL-AGENDA-ESEB-2025.pdf}{eseb.org} \\ Measuring morphological changes in unit of nucleotide changes to disentangle neutral
and adaptive processes.}
\cventry{July 26–30, 2024}{Evolution, Third Joint Congress on Evolutionary Biology}{Montreal.}{}{}{\href{https://youtu.be/lXyhTDWSVfg?feature=shared&t=4240}{youtu.be/lXyhTDWSVfg?feature=shared\&t=4240} \\ Detecting diversifying selection for a trait from within and between-species genotypes and phenotypes.}
\cventry{July 12–16, 2024}{Intelligent Systems for Molecular Biology}{Montreal.}{}{}{\href{https://www.iscb.org/ismb2024/programme-schedule/scientific-programme/vari}{iscb.org/ismb2024} \\ A phylogenetic mutation-selection model predicts fitness effects of mutations in extant mammals.}
\cventry{February 16–17, 2023}{Biology 2023}{Geneve.}{}{}{\href{https://biology23.unige.ch/wp-content/uploads/2023/02/Biology23_Conference_Booklet.pdf}{biology23.unige.ch} \\ Up to 25\% of beneficial mutations in protein sequences are not adaptive innovations in mammals.}
\cventry{January 23–25, 2023}{Interdisciplinary Approaches for Molecular Evolution}{Grenoble.}{}{}{\href{https://alphy-aiem-2023.sciencesconf.org/data/pages/Alphy_AIEM_Janv2023_UGA_Abstracts_booklet.pdf}{alphy-aiem-2023.sciencesconf.org} \\ Up to 25\% of beneficial mutations in protein sequences are not adaptive innovations in mammals.}
\cventry{August 14–19, 2022}{European Society for Evolutionary Biology}{Prague.}{}{}{Empirical evidence for positive selection that is not adaptive evolution.}
\cventry{June 26–30, 2022}{Mathematical \& Computational Evolutionary Biology}{Chateau d’Oex.}{}{}{\href{https://mceb2022.sciencesconf.org/data/pages/MCEB2022_program_3.pdf}{mceb2022.sciencesconf.org} \\Empirical evidence for positive selection that is not adaptive evolution.}
\cventry{June 29, 2020}{Society for Molecular Biology \& Evolution}{Quebec [cancelled].}{}{}{Reconstructing changes in population size at the phylogenetic scale from the pattern of substitutions.}
\cventry{July 21-25, 2019}{Society for Molecular Biology \& Evolution}{Manchester.}{}{}{\href{https://smbe.org/smbe/SMBE2019Meeting/www.smbe2019.org/media/1058/smbe-2019-abstract-book.pdf}{smbe.org/smbe/SMBE2019Meeting/www.smbe2019.org/index.html} \\Inferring fluctuating population size and selection with phylogenetic codon models.}
\cventry{July 16-19, 2019}{Mathematical Models in Ecology \& Evolution}{Lyon.}{}{}{\href{https://mmee2019lyon.sciencesconf.org/data/pages/program_MMEE_full_v5.pdf}{mmee2019lyon.sciencesconf.org} \\Inferring fluctuating population size and selection with phylogenetic codon models.}
\cventry{August 21, 2018}{Evolution, Second Joint Congress on Evolutionary Biology}{Montpellier.}{}{}{\href{https://programme.europa-organisation.com/slides/programme_jointCongressEvolBiology-2018/webconf/862_21082018_1120_rabelais_Thibault_Latrille_697/index.html}{programme.europa-organisation.com/slides/programme\_jointCongressEvolBiology-2018/webconf/862\_21082018\_1120\_rabelais\_Thibault\_Latrille\_697/index.html}\\The Red-Queen model of recombination hotspots evolution.}
\cventry{November 8-9, 2017}{Interdisciplinary Approaches for Molecular Evolution}{Lyon.}{}{}{\href{https://project.inria.fr/aiem2017/files/2017/07/AIEM_2017-Abstract_07112017.pdf}{project.inria.fr/aiem2017} \\ Application of mean-field theory in red-queen dynamics.}
\cventry{November 2-3, 2016}{Blend Web Mix}{Lyon.}{}{}{\href{https://www.youtube.com/watch?v=oH_-9-z-o_c}{youtube.com/watch?v=oH\_-9-z-o\_c} \\Le web \& internet ont-ils changé notre pratique de la science ? Mariages, infidélités et divorces entre la science et le numérique.}
\cventry{June 6-7, 2016}{Modèles en Ecologie Evolutive}{Montpellier.}{}{}{La reine rouge au royaume des recombinaisons.}
\cventry{June 3-4, 2014}{Modèles en Ecologie Evolutive}{Montpellier.}{}{}{Robust estimation of phylogenetic diversity: steer clear of rare species.}
\cventry{May 27-28, 2014}{Bioinformatics for Environmental Genomics}{Lyon.}{}{}{\href{https://ge-lyon2014.sciencesconf.org/conference/ge-lyon2014/Abstracts.pdf}{ge-lyon2014.sciencesconf.org} \\ Robust estimation of phylogenetic diversity: steer clear of rare species.}
\end{document}